\chapter{Compacte operatoren en de spectraalstelling}
\label{ch:b3:spectral}

Diagonaliseren is het kroonjuweel van de eindigdimensionale
lineaire algebra: een symmetrische matrix heeft een orthonormale
basis van eigenvectoren. In oneindige dimensie faalt dat voor
begrensde \omterm{def:b3:spectral:selfadjoint}{zelftoegevoegde operatoren} in het algemeen ---
vermenigvuldiging met $x$ op $L^2(\intcc01)$ heeft
\emph{helemaal geen} eigenwaarden (\cref{exo:b3:spectral:6}) ---
maar het overleeft, in volmaakte vorm, voor de operatoren die
\emph{bijna eindigdimensionaal} zijn: de \omterm{def:b3:topology:compact}{compacte}. De
spectraalstelling voor \omterm{def:b3:topology:compact}{compacte} \omterm{def:b3:spectral:selfadjoint}{zelftoegevoegde operatoren} is de
meest gebruikte stelling van de toegepaste functionaalanalyse: ze
diagonaliseert integraalvergelijkingen, drijft het \omterm{thm:b3:spectral:fredholm}{alternatief
van Fredholm} aan en lost (in de weekendopgave) de trillende snaar
op, waarbij ze de sinusbasis van de fourieranalyse uit pure
operatortheorie voortbrengt --- met Eulers $\zeta(2) =
\frac{\pi^2}6$ als afscheidsgeschenk uit een spoorformule.
Overal is $H$ een \omterm{def:b3:hilbert:inner}{hilbertruimte} over $\C$ (of $\R$; de
uitspraken passen zich aan), en zijn de operatoren begrensd.

\section{Compacte operatoren}

\begin{definition}\label{def:b3:spectral:compact}
$T \in \mathcal L(E, F)$ ($E, F$ banach) heet
\emph{compact}\index{compacte operator} als het beeld $T(B)$ van
de eenheidsbal relatief compact is in $F$ --- gelijkwaardig: als
elke begrensde rij $(x_n)$ een deelrij heeft waarvoor
$(Tx_{n_k})$ convergeert. Operatoren van eindige rang zijn
compact (begrensde verzamelingen in eindige dimensie); de
identiteit van een oneindigdimensionale ruimte is dat nooit (de
stelling van Riesz, volume van bachelorjaar 2).
\end{definition}

\begin{proposition}\label{prop:b3:spectral:ideal}
De \omterm{def:b3:spectral:compact}{compacte operatoren} $\mathcal K(E, F)$ vormen een gesloten
deelruimte van $\mathcal L(E, F)$ en een tweezijdig \omterm{def:b3:rings:ideal}{ideaal}: is
$S$ \omterm{def:b3:spectral:compact}{compact}, dan zijn $AS$ en $SB$ \omterm{def:b3:spectral:compact}{compact} voor begrensde $A, B$.
Bovendien is in een \omterm{def:b3:hilbert:inner}{hilbertruimte} elke \omterm{def:b3:spectral:compact}{compacte operator} een
limiet in norm van operatoren van eindige rang.
\end{proposition}

\begin{proof}
Deelruimte: duidelijk uit de karakterisering met rijen. \omterm{def:b3:rings:ideal}{Ideaal}:
begrensde afbeeldingen sturen convergente rijen naar convergente
en begrensde naar begrensde. Geslotenheid: zij $T_n \to T$ met
$T_n$ \omterm{def:b3:spectral:compact}{compact}, en $(x_k)$ begrensd door $1$; een
diagonaalextractie maakt $(T_nx_{k_j})_j$ convergent voor elke
$n$; dan is $(Tx_{k_j})$ cauchy, want
\[
\norm{Tx_{k_j} - Tx_{k_l}} \leq 2\vertiii{T - T_n} +
\norm{T_nx_{k_j} - T_nx_{k_l}} ,
\]
waarbij je eerst $n$ kiest en dan de indices. Benadering in
\omterm{def:b3:hilbert:inner}{hilbertruimten}: zij $T$ \omterm{def:b3:spectral:compact}{compact} en $K = \overline{T(B)}$
\omterm{def:b3:spectral:compact}{compact}; overdek $K$ bij gegeven $\varepsilon$ met eindig veel
ballen $B(y_i, \varepsilon)$ en zij $P$ de orthogonale projectie
op $V = \operatorname{Vect}(y_1, \dots, y_m)$ (gesloten:
eindigdimensionaal). Dan heeft $PT$ eindige rang, en voor $\norm
x \leq 1$ geldt, met een $y_i$ waarvoor $\norm{Tx - y_i} <
\varepsilon$,
\[
\norm{Tx - PTx} \leq \norm{Tx - y_i} + \norm{P(y_i - Tx)}
\leq 2\varepsilon
\]
($y_i = Py_i$; $\vertiii P \leq 1$): dus $\vertiii{T - PT} \leq
2\varepsilon$.
\end{proof}

\begin{example}\label{ex:b3:spectral:examples}
(a) Diagonaaloperatoren op $\ell^2$: $T(x_n) = (\lambda_nx_n)$ is
\omterm{def:b3:spectral:compact}{compact} dan en slechts dan als $\lambda_n \to 0$
(\cref{exo:b3:spectral:2}).
(b) Kernoperatoren op $\mathcal C(\intcc01)$: \omterm{def:b3:spectral:compact}{compact} volgens
Ascoli (\cref{exo:b3:complete:7}).
(c) \emph{Hilbert-schmidtoperatoren}\index{hilbert-schmidtoperator}:
voor $k \in L^2(\intcc01^2)$ definieert
\[
(T_kf)(x) = \int_0^1k(x, y)\,f(y)\,\dd y
\]
een \omterm{def:b3:spectral:compact}{compacte operator} op $L^2(\intcc01)$ met $\vertiii{T_k} \leq
\norm k_{L^2}$ (\cref{exo:b3:spectral:4}: de basisontwikkeling
van $k$ afknotten toont $T_k$ als limiet van operatoren van
eindige rang).
\end{example}

\section{Zelftoegevoegde operatoren}

\begin{definition}\label{def:b3:spectral:selfadjoint}
$T \in \mathcal L(H)$ heet
\emph{zelftoegevoegd}\index{zelftoegevoegde operator} als $T =
T^*$ (\cref{exo:b3:hilbert:8}), dat wil zeggen $\langle Tx,
y\rangle = \langle x, Ty\rangle$ voor alle $x, y$. Dan is
$\langle x, Tx\rangle \in \R$ voor elke $x$ (gelijk aan haar
toegevoegde).
\end{definition}

\begin{proposition}\label{prop:b3:spectral:sanorm}
Voor \omterm{def:b3:spectral:selfadjoint}{zelftoegevoegde} $T$ geldt\index{numerieke straal}
\[
\vertiii T = \sup_{\norm x \leq 1}\ \abs{\langle x,
Tx\rangle} .
\]
De eigenwaarden van $T$ zijn reëel, en eigenvectoren bij
verschillende eigenwaarden staan loodrecht op elkaar.
\end{proposition}

\begin{proof}
Zij $M$ het supremum; $M \leq \vertiii T$ volgens
Cauchy--Schwarz. Omgekeerd geeft de identiteit van
polarisatietype
\[
\langle x{+}y, T(x{+}y)\rangle - \langle x{-}y,
T(x{-}y)\rangle = 4\operatorname{Re}\langle y, Tx\rangle
\]
(ontwikkel; de kruistermen $\langle y, Tx\rangle + \langle x,
Ty\rangle = 2\operatorname{Re}\langle y, Tx\rangle$ wegens de
\omterm{def:b3:spectral:selfadjoint}{zelftoegevoegdheid}), samen met de parallellogramwet,
\[
4\operatorname{Re}\langle y, Tx\rangle \leq
M\bigl(\norm{x{+}y}^2 + \norm{x{-}y}^2\bigr) = 2M\bigl(\norm
x^2 + \norm y^2\bigr).
\]
Neem voor $\norm x = 1$ met $Tx \neq 0$ $y = Tx/\norm{Tx}$:
$4\norm{Tx} \leq 4M$. Dus $\vertiii T \leq M$. Eigenwaarden: uit
$Tx = \lambda x$ met $x \ne 0$ volgt $\lambda\norm x^2 = \langle
x, Tx\rangle \in \R$. Orthogonaliteit: $\lambda\langle x,
y\rangle = \langle Tx, y\rangle = \langle x, Ty\rangle =
\mu\langle x, y\rangle$ met $\lambda \neq \mu$ reëel.
\end{proof}

\section{De spectraalstelling}

\begin{lemma}[Bestaan van een extreme eigenwaarde]
\label{lem:b3:spectral:existence}
Zij $T \neq 0$ \omterm{def:b3:spectral:compact}{compact} en \omterm{def:b3:spectral:selfadjoint}{zelftoegevoegd}. Dan is $\vertiii T$ of
$-\vertiii T$ een eigenwaarde van $T$.
\end{lemma}

\begin{proof}
Kies volgens \cref{prop:b3:spectral:sanorm} eenheidsvectoren
$x_n$ met $\langle x_n, Tx_n\rangle \to \mu$, waarbij $\abs\mu =
\vertiii T > 0$ (ga over op een deelrij om het teken vast te
leggen). Dan is
\[
\norm{Tx_n - \mu x_n}^2
= \norm{Tx_n}^2 - 2\mu\langle x_n, Tx_n\rangle + \mu^2
\leq \vertiii T^2 - 2\mu\langle x_n, Tx_n\rangle + \mu^2
\longrightarrow 2\mu^2 - 2\mu\cdot\mu = 0 .
\]
Wegens de \omterm{def:b3:topology:compact}{compactheid} convergeert een deelrij $Tx_{n_k} \to y$;
dan is $\mu x_{n_k} = Tx_{n_k} - (Tx_{n_k} - \mu x_{n_k}) \to y$,
dus $x_{n_k} \to x = y/\mu$, een eenheidsvector, en de
\omterm{def:b3:topology:continuity}{continuïteit} geeft $Tx = \mu x$.
\end{proof}

\begin{theorem}[Spectraalstelling voor compacte zelftoegevoegde
operatoren]\label{thm:b3:spectral:spectral}
Zij $T$ een \omterm{def:b3:topology:compact}{compacte} \omterm{def:b3:spectral:selfadjoint}{zelftoegevoegde operator} op een
\omterm{def:b3:hilbert:inner}{hilbertruimte} $H$.\index{spectraalstelling (compact
zelftoegevoegd)}
\begin{enumerate}
  \item $H$ laat een orthonormaal stelsel $(e_n)_{n \in N}$ ($N$
        eindig of aftelbaar) van eigenvectoren van $T$ toe, met
        reële eigenwaarden $(\lambda_n) \neq 0$, zodat
        \[
        Tx = \sum_{n\in N}\lambda_n\,\langle e_n, x\rangle\,
        e_n \qquad (x \in H),
        \]
        en $H = \ker T \,\oplus^\perp\,
        \overline{\operatorname{Vect}}(e_n : n \in N)$.
  \item Is $N$ oneindig, dan is $\lambda_n \to 0$; voor elke
        $\delta > 0$ hebben slechts eindig veel $n$
        $\abs{\lambda_n} \geq \delta$, en elke eigenruimte
        $\ker(T - \lambda)$ met $\lambda \neq 0$ is
        eindigdimensionaal.
  \item $(e_n)$ aanvullen met een orthonormale basis van $\ker T$
        levert, wanneer $H$ separabel is, een orthonormale basis
        van $H$ die uit eigenvectoren bestaat: $T$ is
        gediagonaliseerd.
\end{enumerate}
\end{theorem}

\begin{proof}
Eerst (2). Hadden oneindig veel orthonormale eigenvectoren $x_k$
$\abs{\lambda_{(k)}} \geq \delta$, dan is $\norm{Tx_k - Tx_l}^2 =
\lambda_{(k)}^2 + \lambda_{(l)}^2 \geq 2\delta^2$
(orthogonaliteit, Pythagoras): geen enkele convergente deelrij
van $(Tx_k)$, in tegenspraak met de \omterm{def:b3:topology:compact}{compactheid} van $T$ op de
begrensde $(x_k)$. Dat begrenst voor elke $\delta$ de totale
multipliciteit van de eigenwaarden buiten $\intoo{-\delta}\delta$
door een eindig getal; de aftelbaarheid en $\lambda_n \to 0$
volgen.

(1) Zij $H_0$ het gesloten opspansel van \emph{alle}
eigenvectoren met eigenwaarde $\neq 0$, geordend (met (2) en
Gram--Schmidt binnen elke eindigdimensionale eigenruimte, en de
orthogonaliteit tussen eigenruimten uit
\cref{prop:b3:spectral:sanorm}) tot een orthonormaal stelsel
$(e_n)$ met eigenwaarden $\lambda_n \neq 0$. $T$ beeldt $H_0$ in
$H_0$ af, en ook $H_0^\perp$ in $H_0^\perp$: voor $y \perp H_0$
en $e$ een eigenvector is $\langle e, Ty\rangle = \langle Te,
y\rangle = \lambda\langle e, y\rangle = 0$. De beperking $T' =
T\restriction_{H_0^\perp}$ is \omterm{def:b3:spectral:compact}{compact} en \omterm{def:b3:spectral:selfadjoint}{zelftoegevoegd} op de
\omterm{def:b3:hilbert:inner}{hilbertruimte} $H_0^\perp$; was $T' \neq 0$, dan zou
\cref{lem:b3:spectral:existence} binnen $H_0^\perp$ een
eigenvector van $T$ met eigenwaarde $\neq 0$ voortbrengen ---
onmogelijk, want zulke vectoren wonen in $H_0$. Dus $T' = 0$:
$H_0^\perp \subseteq \ker T$. Omgekeerd staat $\ker T$ loodrecht
op elke $e_n$ ($\langle e_n, z\rangle = \frac1{\lambda_n}\langle
Te_n, z\rangle = \frac1{\lambda_n}\langle e_n, Tz\rangle = 0$):
$\ker T \subseteq H_0^\perp$, waaruit $\ker T = H_0^\perp$ en de
orthogonale ontbinding volgen. De ontwikkeling: voor $x = z +
\sum_nc_ne_n$ ($z \in \ker T$, $c_n = \langle e_n, x\rangle$;
\cref{thm:b3:hilbert:parseval}(1) op $H_0$) geeft de \omterm{def:b3:topology:continuity}{continuïteit}
van $T$ dat $Tx = \sum_nc_n\lambda_ne_n$.

(3) $\ker T$, een gesloten deelruimte van een separabele ruimte,
is separabel: ze heeft een orthonormale basis
(\cref{prop:b3:hilbert:gramschmidt}); de vereniging is wegens de
ontbinding in (1) een orthonormale basis van $H$.
\end{proof}

\begin{theorem}[Alternatief van Fredholm]
\label{thm:b3:spectral:fredholm}
Zij $T$ \omterm{def:b3:spectral:compact}{compact} en \omterm{def:b3:spectral:selfadjoint}{zelftoegevoegd} en $\lambda \in
\R\setminus\{0\}$.\index{alternatief van Fredholm}
\begin{enumerate}
  \item Is $\lambda$ geen eigenwaarde, dan is $T - \lambda I$
        bijectief met begrensde inverse: voor elke $f$ heeft de
        vergelijking $Tx - \lambda x = f$ precies één oplossing,
        die \omterm{def:b3:topology:continuity}{continu} van $f$ afhangt.
  \item Is $\lambda$ een eigenwaarde, dan is $Tx - \lambda x = f$
        \omterm{def:b3:groups:derived}{oplosbaar} dan en slechts dan als $f \perp \ker(T -
        \lambda I)$, en de oplossing is uniek op die
        (eindigdimensionale) kern na.
\end{enumerate}
\end{theorem}

\begin{proof}
Ontbind $x = z + \sum c_ne_n$ en $f = w + \sum d_ne_n$ langs
\cref{thm:b3:spectral:spectral} ($z, w \in \ker T$). De
vergelijking luidt
\[
-\lambda z = w, \qquad (\lambda_n - \lambda)\,c_n = d_n\
(n \in N).
\]
(1) $\lambda \notin \{\lambda_n\}\cup\{0\}$: volgens (2) van de
spectraalstelling is $\inf_n\abs{\lambda_n - \lambda} = \delta >
0$ (de eigenwaarden hopen zich enkel op in $0 \neq \lambda$). Los
op: $z = -w/\lambda$ en $c_n = d_n/(\lambda_n - \lambda)$, met
$\sum\abs{c_n}^2 \leq \delta^{-2}\sum\abs{d_n}^2$: één unieke
oplossing met $\norm x \leq C\norm f$.
(2) $\lambda = \lambda_{n}$ voor $n$ in een eindige verzameling
$F$: de \omterm{def:b3:groups:derived}{oplosbaarheid} van $(\lambda_n - \lambda)c_n = d_n$ voor
$n \in F$ vereist $d_n = 0$, dat wil zeggen $f \perp e_n$ ($n \in
F$), oftewel $f \perp \ker(T - \lambda I)$; de $c_n$ met $n \in
F$ zijn dan vrij.
\end{proof}

\begin{example}\label{ex:b3:spectral:minxy}
Stel op $L^2(\intcc01)$ $Tf(x) = \int_0^1\min(x, y)f(y)\dd y$:
een \omterm{ex:b3:spectral:examples}{hilbert-schmidtoperator} met reële symmetrische kern, dus
\omterm{def:b3:spectral:compact}{compact} en \omterm{def:b3:spectral:selfadjoint}{zelftoegevoegd}. Het oplossen van $Tf = \lambda f$: de
betrekking $\bigl(Tf\bigr)(x) = \int_0^xyf(y)\dd y +
x\int_x^1f(y)\dd y$ toont dat $u = Tf$ voldoet aan $u'' = -f$
(twee keer differentiëren, geoorloofd voor \omterm{def:b3:topology:continuity}{continue} $f$, en $Tf$
is \omterm{def:b3:topology:continuity}{continu} voor $f \in L^2$: gedomineerde convergentie), met
$u(0) = 0$ en $u'(1) = 0$. Eigenfuncties lossen dus $\lambda u''
= -u$ op met $u(0) = 0$ en $u'(1) = 0$:
\[
u_n(x) = \sin\Bigl(\bigl(n + \tfrac12\bigr)\pi x\Bigr),
\qquad
\lambda_n = \frac{1}{\bigl(n + \frac12\bigr)^2\pi^2}
\quad (n \geq 0),
\]
en de spectraalstelling beweert --- zonder enige fouriertheorie
--- dat deze sinussen na normering een orthonormale basis van
$L^2(\intcc01)$ vormen (de kern van $T$ is $0$: $Tf = 0$ dwingt
via de twee differentiaties $f = 0$ b.o.\ af). De weekendopgave
doorloopt dezelfde ideeënkring voor de trillende snaar en haalt
$\zeta(2)$ uit het spoor.
\end{example}

\begin{method}\label{met:b3:spectral:method}
Bij een integraal- of differentiaalvergelijking: (1) herschrijf
haar als $(I - \lambda K)u = f$ of $Ku = \lambda u$ met $K$ een
integraaloperator; (2) ga na dat $K$ \omterm{def:b3:spectral:compact}{compact} is
(hilbert-schmidtkern, of Ascoli) en indien mogelijk
\omterm{def:b3:spectral:selfadjoint}{zelftoegevoegd} (reële symmetrische kern); (3) diagonaliseer met
de spectraalstelling, of roep het \omterm{thm:b3:spectral:fredholm}{alternatief van Fredholm} in
voor de \omterm{def:b3:groups:derived}{oplosbaarheid}; (4) lees bestaan, uniciteit, stabiliteit
en reeksformules voor de oplossingen in de eigenbasis af.
Differentiaaloperatoren zijn onbegrensd, maar hun \emph{inversen}
(greenoperatoren) zijn \omterm{def:b3:spectral:compact}{compact}: inverteer altijd eerst.
\end{method}

\section{Oefeningen}

\begin{exercise}[$\star$]\label{exo:b3:spectral:1}
(a) Toon aan dat een begrensde operator met eindigdimensionaal
beeld \omterm{def:b3:spectral:compact}{compact} is.
(b) Toon aan dat de identiteit van een genormeerde ruimte \omterm{def:b3:spectral:compact}{compact}
is dan en slechts dan als de dimensie eindig is (Riesz, volume
van bachelorjaar 2). Leid af dat een \omterm{def:b3:spectral:compact}{compacte operator} op een
oneindigdimensionale ruimte nooit inverteerbaar is met begrensde
inverse.
\end{exercise}

\begin{exercise}[$\star$]\label{exo:b3:spectral:2}
Zij $T(x_1, x_2, \dots) = (\lambda_1x_1, \lambda_2x_2, \dots)$ op
$\ell^2$, met $(\lambda_n)$ begrensd.
(a) Toon aan dat $\vertiii T = \sup\abs{\lambda_n}$.
(b) Toon aan dat $T$ \omterm{def:b3:spectral:compact}{compact} is dan en slechts dan als $\lambda_n
\to 0$. \emph{(Voor $\Leftarrow$ afknotten; voor $\Rightarrow$
testen op $(e_n)$.)}
(c) Wanneer is $T$ \omterm{def:b3:spectral:selfadjoint}{zelftoegevoegd}? Ga in dat geval de
spectraalstelling met het blote oog na.
\end{exercise}

\begin{exercise}[$\star\star$]\label{exo:b3:spectral:3}
Geef de details van de ideaaleigenschap
(\cref{prop:b3:spectral:ideal}): is $S$ \omterm{def:b3:spectral:compact}{compact} en zijn $A, B$
begrensd, dan is $ASB$ \omterm{def:b3:spectral:compact}{compact}. Leid af dat als $ST = TS = I$
voor een begrensde $S$ en $\dim H = \infty$, $T$ niet \omterm{def:b3:spectral:compact}{compact} is
--- en verzoen dat met \cref{exo:b3:spectral:1}(b).
\end{exercise}

\begin{exercise}[$\star\star$]\label{exo:b3:spectral:4}
(Hilbert--Schmidt) Zij $k \in L^2(\intcc01^2)$ en $(e_n)$ een
\omterm{def:b3:hilbert:onb}{hilbertbasis} van $L^2(\intcc01)$.
(a) Toon aan dat $\vertiii{T_k} \leq \norm k_{L^2}$
\emph{(Cauchy--Schwarz in de veranderlijke $y$, daarna
Tonelli)}.
(b) Ontwikkel $k(x,y) = \sum_{m,n}c_{mn}e_m(x)\overline{e_n(y)}$
in $L^2$ van het vierkant (rechtvaardig dat de producten daar een
\omterm{def:b3:hilbert:onb}{hilbertbasis} vormen), en toon aan dat de afgeknotte som
operatoren van eindige rang geeft die in \omterm{def:b3:banach:operator}{operatornorm} naar $T_k$
convergeren: dus is $T_k$ \omterm{def:b3:spectral:compact}{compact}.
\end{exercise}

\begin{exercise}[$\star\star$]\label{exo:b3:spectral:5}
Zij $T$ \omterm{def:b3:spectral:selfadjoint}{zelftoegevoegd} met $\langle x, Tx\rangle \geq 0$ voor alle
$x$ (een \emph{positieve} operator).
(a) Toon aan dat de eigenwaarden $\geq 0$ zijn en dat $\vertiii T
= \sup_{\norm x\leq1}\langle x, Tx\rangle$.
(b) Bewijs de veralgemeende ongelijkheid van Cauchy--Schwarz
$\abs{\langle x, Ty\rangle}^2 \leq \langle x, Tx\rangle\langle y,
Ty\rangle$.
\end{exercise}

\begin{exercise}[$\star\star$]\label{exo:b3:spectral:6}
Stel op $L^2(\intcc01)$ $(Mf)(x) = x\,f(x)$.
(a) Toon aan dat $M$ begrensd en \omterm{def:b3:spectral:selfadjoint}{zelftoegevoegd} is met $\vertiii
M = 1$, maar \emph{geen} eigenwaarden heeft.
(b) Toon aan dat $M$ niet \omterm{def:b3:spectral:compact}{compact} is \emph{(geef een begrensde
rij waarvan het beeld geen convergente deelrij heeft, bijvoorbeeld
genormeerde indicatoren van krimpende intervallen nabij $1$ ---
of roep de spectraalstelling in)}.
(c) Waar breekt het bewijs van \cref{lem:b3:spectral:existence}
voor $M$?
\end{exercise}

\begin{exercise}[$\star\star$]\label{exo:b3:spectral:7}
(Volterra) Stel op $L^2(\intcc01)$ $Vf(x) = \int_0^xf(y)\dd y$.
(a) Toon aan dat $V$ \omterm{def:b3:spectral:compact}{compact} is (\omterm{ex:b3:spectral:examples}{hilbert-schmidtoperator} met kern
$\mathbf 1_{y < x}$) maar niet \omterm{def:b3:spectral:selfadjoint}{zelftoegevoegd}; bereken $V^*$.
(b) Toon aan dat $V$ geen eigenwaarde $\neq 0$ heeft. \emph{(Uit
$Vf = \lambda f$: $f$ heeft een \omterm{def:b3:topology:continuity}{continue} vertegenwoordiger, is
dan $\mathcal C^1$, en lost $\lambda f' = f$ met $f(0) = 0$
op.)}
(c) Besluit dat \omterm{def:b3:topology:compact}{compactheid} alleen geen eigenvectoren oplevert:
de \omterm{def:b3:spectral:selfadjoint}{zelftoegevoegdheid} in \cref{thm:b3:spectral:spectral} is
essentieel.
\end{exercise}

\begin{exercise}[$\star\star\star$]\label{exo:b3:spectral:8}
(Courant--Fischer) Zij $T$ \omterm{def:b3:spectral:compact}{compact}, \omterm{def:b3:spectral:selfadjoint}{zelftoegevoegd} en
\emph{positief}, met eigenwaarden $\mu_1 \geq \mu_2 \geq \dots >
0$ (herhaald naar multipliciteit, met eigenvectoren $e_1, e_2,
\dots$). Toon aan:
\[
\mu_{k} = \max_{\substack{V \subseteq H \\ \dim V = k}}\
\min_{\substack{x \in V\\ \norm x = 1}}\ \langle x, Tx\rangle
= \min_{\substack{W \subseteq H\\ \operatorname{codim}W =
k-1}}\ \max_{\substack{x\in W\\ \norm x = 1}}\ \langle x,
Tx\rangle .
\]
\emph{(Test $V = \operatorname{Vect}(e_1,\dots,e_k)$; snijd voor
de bovengrens een willekeurige $V$ met ruimten van het type
$\operatorname{Vect}(e_k, e_{k+1}, \dots)$: het tellen van
dimensies dwingt een doorsnede $\neq 0$ af.)} Leid af dat de
eigenwaarden monotoon van $T$ afhangen ($T \leq S \Rightarrow
\mu_k(T) \leq \mu_k(S)$).
\end{exercise}

\begin{exercise}[$\star\star$]\label{exo:b3:spectral:9}
Voor welke $\lambda \in \R$ heeft de integraalvergelijking
\[
f(x) - \lambda\int_0^1\min(x,y)\,f(y)\,\dd y = g(x)
\]
met \cref{thm:b3:spectral:fredholm} voor $Tf(x) =
\int_0^1\min(x,y)f(y)\dd y$ (\cref{ex:b3:spectral:minxy}) een
unieke oplossing $f \in L^2$ voor elke $g \in L^2$? Wat gebeurt
er bij de uitzonderlijke waarden?
\end{exercise}

\begin{exercise}[$\star\star\star$]\label{exo:b3:spectral:10}
Zij $S$ de verschuiving op $\ell^2$ (\cref{exo:b3:banach:1}).
(a) Toon aan dat $S$ geen eigenwaarden heeft, terwijl elke
$\lambda$ met $\abs\lambda < 1$ een eigenwaarde van $S^*$ is
(zoek de eigenvectoren expliciet: meetkundige rijen).
(b) Noch $S$ noch $S^*$ is \omterm{def:b3:spectral:compact}{compact}: ga dat na door te testen op
$(e_n)$ in de trant van \cref{exo:b3:spectral:2}.
(c) Becommentarieer: voor \omterm{def:b3:spectral:selfadjoint}{niet-zelftoegevoegde}, \omterm{def:b3:topology:compact}{niet-compacte}
operatoren kan het eigenwaardenlandschap alles zijn, van leeg tot
een volle schijf --- het begrip dat overleeft is het
\emph{spectrum}, dat in een latere cursus wordt bestudeerd.
\end{exercise}

\begin{exercise}[$\star\star$]\label{exo:b3:spectral:11}
(Vierkantswortels) Zij $T$ \omterm{def:b3:spectral:compact}{compact}, \omterm{def:b3:spectral:selfadjoint}{zelftoegevoegd} en positief
($\langle x, Tx\rangle \geq 0$) op een \omterm{def:b3:hilbert:inner}{hilbertruimte} $H$, met
spectrale ontbinding $Tx = \sum_n\mu_n\langle e_n, x\rangle e_n$
($\mu_n > 0$).
(a) Definieer $Sx = \sum_n\sqrt{\mu_n}\,\langle e_n, x\rangle
e_n$; toon aan dat $S$ \omterm{def:b3:spectral:compact}{compact}, \omterm{def:b3:spectral:selfadjoint}{zelftoegevoegd} en positief is met
$S^2 = T$.
(b) Bewijs de uniciteit: elke \omterm{def:b3:topology:compact}{compacte} positieve \omterm{def:b3:spectral:selfadjoint}{zelftoegevoegde}
$R$ met $R^2 = T$ behoudt de eigenruimten van $T$ \emph{($RT =
R^3 = TR$: $R$ commuteert met $T$, dus $R(\ker(T - \mu))
\subseteq \ker(T - \mu)$)}, en op $\ker(T - \mu)$ is $R$ een
positieve operator die op een eindigdimensionale ruimte kwadrateert
tot $\mu\,\mathrm{id}$: diagonaliseer haar daar en besluit dat $R
= \sqrt\mu\,\mathrm{id}$ op elke eigenruimte, dus $R = S$.
(c) Bereken $\sqrt G$ voor de snaaroperator $G$ van
\cref{pb:b3:spectral:1}: welke kern heeft op de sinusbasis de
eigenwaarden $\frac1{n\pi}$? (Druk $\sqrt G$ uit als
$L^2$-limiet van kernen; een gesloten vorm is niet vereist.)
\end{exercise}

\begin{exercise}[$\star\star\star$]\label{exo:b3:spectral:12}
(Singulierewaardenontbinding) Zij $T \in \mathcal L(H)$
\emph{\omterm{def:b3:spectral:compact}{compact}}, niet noodzakelijk \omterm{def:b3:spectral:selfadjoint}{zelftoegevoegd}.
(a) Toon aan dat $T^*T$ \omterm{def:b3:spectral:compact}{compact}, \omterm{def:b3:spectral:selfadjoint}{zelftoegevoegd} en positief is;
zij $(e_n)$ een orthonormale familie eigenvectoren met $T^*Te_n =
s_n^2e_n$ en $s_n > 0$ (de \emph{singuliere waarden}), aangevuld
met $\ker(T^*T) = \ker T$ (bewijs die gelijkheid).
(b) Stel $f_n = \frac{Te_n}{s_n}$; toon aan dat $(f_n)$
orthonormaal is, en vestig de \emph{singulierewaardenontbinding}
\[
Tx = \sum_n s_n\,\langle e_n, x\rangle\,f_n
\qquad (x \in H),
\]
met convergentie in $H$.
(c) Leid af: $\vertiii T = \max_ns_n$; $T$ is een limiet in norm
van operatoren van eindige rang (wat de omkering in
\cref{prop:b3:spectral:ideal} voor \omterm{def:b3:hilbert:inner}{hilbertruimten} opnieuw
bewijst); en leg voor de volterra-operator $V$ van
\cref{exo:b3:spectral:7}, die geen eigenwaarden heeft, uit
waarom de ontbinding niettemin bestaat en wat haar ingrediënten
zijn (identificeer $V^*V$ als een kernoperator van het snaartype
--- haar eigenwaarden expliciet berekenen is het terrein van
\cref{exo:b3:spectral:9}).
\end{exercise}

\section{Probleem: de trillende snaar en
\texorpdfstring{$\zeta(2)$}{zeta(2)}}

\begin{problem}[{Weekendopgave --- de operator van Green, de
sinusbasis en een spoorformule}]\label{pb:b3:spectral:1}
We lossen het eigenwaardeprobleem van de trillende snaar met
vaste uiteinden --- $-u'' = \nu u$, $u(0) = u(1) = 0$ --- met
operatortheorie op, verkrijgen de orthonormale sinusbasis zonder
enige fourierberekening, en bepalen $\zeta(2)$ door twee
uitdrukkingen voor het \emph{spoor} van de greenoperator te
vergelijken. Definieer op $L^2(\intcc01)$
\[
(Gf)(x) = \int_0^1 g(x,y)\,f(y)\,\dd y,
\qquad
g(x, y) = \min(x,y)\,\bigl(1 - \max(x,y)\bigr).
\]

\medskip
\noindent\textbf{Deel I --- De greenoperator.}
\begin{enumerate}
  \item Toon aan dat $g$ \omterm{def:b3:topology:continuity}{continu} en symmetrisch is met $0 \leq g
        \leq \frac14$, en dat $G$ \omterm{def:b3:spectral:compact}{compact} en \omterm{def:b3:spectral:selfadjoint}{zelftoegevoegd} is
        (\cref{ex:b3:spectral:examples}(c)).
  \item Toon voor \omterm{def:b3:topology:continuity}{continue} $f$ aan dat $u = Gf$ $\mathcal C^2$ is
        met
        \[
        -u'' = f, \qquad u(0) = u(1) = 0
        \]
        \emph{(schrijf $u(x) = (1-x)\int_0^xyf(y)\dd y +
        x\int_x^1(1-y)f(y)\dd y$ en differentieer tweemaal)}.
        Omgekeerd is, als $u \in \mathcal C^2$ met $u(0) = u(1)
        = 0$, $G(-u'') = u$: $G$ inverteert de snaaroperator.
  \item Toon aan dat $\ker G = \{0\}$ \emph{(is $Gf = 0$ met $f
        \in L^2$: test tegen \omterm{def:b3:topology:continuity}{continue} $\varphi$, breng $G$ met de
        symmetrie en Fubini over op $\varphi$, en gebruik het
        hoofdlemma \cref{cor:b3:lp:fundlemma} --- of strijk
        glad)}, en dat $G$ een positieve operator is: $\langle f,
        Gf\rangle \geq 0$. \emph{(Voor \omterm{def:b3:topology:continuity}{continue} $f$: $\langle f,
        Gf\rangle = \int_0^1 (u')^2$ met $u = Gf$, met partiële
        integratie; besluit met \omterm{ex:b3:lebesgue:gamma}{dichtheid}.)}
\end{enumerate}

\medskip
\noindent\textbf{Deel II --- Diagonaliseren: de sinusbasis.}
\begin{enumerate}[resume]
  \item Toon aan dat de eigenfuncties van $G$ bij eigenwaarde
        $\lambda \ne 0$ op scalairen na de oplossingen zijn van
        $-\lambda u'' = u$ met $u(0) = u(1) = 0$ \emph{(een
        eigenfunctie heeft een \omterm{def:b3:topology:continuity}{continue} vertegenwoordiger ---
        $Gf$ is \omterm{def:b3:topology:continuity}{continu} voor $f \in L^2$, waarom? --- en is dus
        $\mathcal C^2$ door vraag 2 als hefboom te gebruiken)}.
  \item Los het randwaardeprobleem op: de eigenwaarden van $G$
        zijn $\lambda_n = \frac1{n^2\pi^2}$ ($n \geq 1$), met
        genormeerde eigenfuncties $e_n(x) = \sqrt2\,\sin(n\pi
        x)$; ga de orthonormaliteit als controle met een
        rechtstreekse integratie na.
  \item Besluit uit \cref{thm:b3:spectral:spectral} en vraag 3
        dat $\bigl(\sqrt2\sin(n\pi x)\bigr)_{n\geq1}$ een
        orthonormale \emph{basis} van $L^2(\intcc01)$ is --- geen
        Stone--Weierstrass, geen fourierreeksen nodig. Ontwikkel
        $f(x) = x(1-x)$ in deze basis en schrijf \omterm{thm:b3:hilbert:parseval}{Parseval} ervoor
        uit.
\end{enumerate}

\medskip
\noindent\textbf{Deel III --- De spoorformule en $\zeta(2)$.}
\begin{enumerate}[resume]
  \item Bewijs de twee identiteiten
        \[
        \langle e_n, Ge_n\rangle = \lambda_n
        \quad\text{en}\quad
        \sum_{n\geq1}\lambda_n = \int_0^1 g(x,x)\,\dd x .
        \]
        Voor de tweede (de \emph{spoorformule}): ontwikkel voor
        vaste $x$ de functie $g(x, \cdot)$ in de basis $(e_n)$
        --- toon aan dat de coëfficiënten $\lambda_ne_n(x)$ zijn,
        zodat $g(x, \cdot) = \sum_n\lambda_ne_n(x)\,e_n$ in
        $L^2$. Hier zijn de sinussen expliciet: ga rechtstreeks
        na dat $\sum_n\lambda_ne_n(x)e_n(y)$ \emph{uniform}
        convergeert op het vierkant (vergelijk met $\sum
        \frac2{n^2\pi^2}$), zodat haar som \omterm{def:b3:topology:continuity}{continu} is en, met
        voor elke $x$ dezelfde $L^2$-ontwikkeling in $y$, overal
        gelijk is aan $g(x,y)$. Stel $y = x$ en integreer term
        voor term.
  \item Bereken $\int_0^1g(x,x)\dd x = \int_0^1x(1-x)\dd x =
        \frac16$, en besluit
        \[
        \sum_{n\geq1}\frac{1}{n^2\pi^2} = \frac16,
        \qquad\text{dat wil zeggen}\qquad
        \boxed{\ \zeta(2) = \frac{\pi^2}{6}\ } :
        \]
        de som van Euler uit een operatorspoor.
  \item Leid $\zeta(2)$ op een derde manier af: pas \omterm{thm:b3:hilbert:parseval}{Parseval} in
        de sinusbasis toe op de constante functie $\mathbf 1$,
        bereken $\sum_{n \text{ oneven}}\frac1{n^2}$ en besluit.
        Vergelijk daarna de mechanismen: in welke zin is het
        spoorargument van de vragen 7 en 8 ``\omterm{thm:b3:hilbert:parseval}{Parseval}, toegepast
        op de hele kern in één keer''?
\end{enumerate}

\medskip
\noindent\textbf{Deel IV --- De snaar trilt.}
\begin{enumerate}[resume]
  \item (Scheiden van veranderlijken, gesynthetiseerd) Definieer
        voor $f \in L^2$
        \[
        u(t, x) = \sum_{n\geq1}\;c_n\,
        \cos(n\pi t)\,\sqrt2\sin(n\pi x),
        \qquad c_n = \langle e_n, f\rangle .
        \]
        Toon aan dat de reeks voor elke $t$ in $L^2(\intcc01)$
        convergeert, dat $t\mapsto u(t, \cdot)$ \omterm{def:b3:topology:continuity}{continu} is met
        waarden in $L^2$, en dat ze voor $f$ in het opspansel van
        eindig veel $e_n$ de golfvergelijking $\partial_t^2u =
        \partial_x^2u$ oplost met $u(0) = f$, $\partial_tu(0) =
        0$ en vaste uiteinden. De eigenwaarden $n^2\pi^2$ zijn de
        gekwadrateerde frequenties: de harmonischen van de snaar
        --- leg de muzikale interpretatie van
        \cref{thm:b3:spectral:spectral} uit in één alinea.
\end{enumerate}

\medskip
\noindent\textbf{Deel V --- Variationele dividenden: de
machtsmethode, de stabiliteit van Weyl en een strikte grens voor
$\pi$.} Zij $A$ een \omterm{def:b3:topology:compact}{compacte} \omterm{def:b3:spectral:selfadjoint}{zelftoegevoegde} \emph{positieve}
operator met eigenwaarden $\mu_1 \geq \mu_2 \geq \cdots > 0$ en
orthonormale eigenvectoren $(u_n)$; $R(x) = \frac{\langle x,
Ax\rangle}{\norm x^2}$. De min-maxformules zijn
\cref{exo:b3:spectral:8}; hier besteden we ze.
\begin{enumerate}[resume]
  \item (Machtsmethode) Schrijf voor $x \neq 0$ $m_p =
        \sum_n\mu_n^p\abs{\langle u_n, x\rangle}^2$. Toon aan dat
        $m_pm_{p+2} \geq m_{p+1}^2$ (Cauchy--Schwarz), leid de
        keten
        \[
        R(x) \;\leq\; \frac{\langle Ax, Ax\rangle}{\langle x,
        Ax\rangle} \;\leq\; R(Ax) \;\leq\; \mu_1
        \]
        af, en bewijs dat $R(A^kx) \to \mu_1$ als $\langle u_1,
        x\rangle \neq 0$: de operator op een willekeurige
        generieke vector itereren berekent de grootste eigenwaarde
        --- de machtsmethode van de numerieke analyse,
        gecertificeerd.
  \item (Stabiliteit van Weyl) Leid voor \omterm{def:b3:topology:compact}{compacte}
        \omterm{def:b3:spectral:selfadjoint}{zelftoegevoegde} positieve $A, B$ uit
        \cref{exo:b3:spectral:8} af dat
        \[
        \abs{\mu_n(A) - \mu_n(B)} \;\leq\; \vertiii{A - B}
        \qquad\text{voor elke } n :
        \]
        het hele spectrum is $1$-lipschitz in de \omterm{def:b3:banach:operator}{operatornorm} ---
        de eigenwaarden van grote symmetrische stelsels kunnen uit
        benaderingen worden berekend met een gegarandeerde fout.
  \item Pas de grens van Rayleigh toe op $G$ met de testfunctie
        $u(x) = x(1-x)$: los $-w'' = u$ met $w(0) = w(1) = 0$ op
        om $Gu = w = \frac{x(1-x)(1 + x - x^2)}{12}$ te krijgen,
        bereken
        \[
        \norm u_2^2 = \frac1{30}, \qquad
        \langle u, Gu\rangle = \frac1{12}\Bigl(\frac1{30} +
        \frac1{140}\Bigr) = \frac{17}{5040},
        \qquad R(u) = \frac{17}{168},
        \]
        en besluit tot de \emph{strikte} grens $\frac1{\pi^2} =
        \lambda_1 \geq \frac{17}{168}$, dat wil zeggen $\pi \leq
        \sqrt{168/17} < 3.1437$.
  \item Eén stap van de keten uit vraag 11, op dezelfde
        testfunctie: bereken met $\int_0^1(x - x^2)^4\dd x =
        \frac1{630}$
        \[
        \norm{Gu}_2^2 = \frac1{144}\Bigl(\frac1{30} +
        \frac2{140} + \frac1{630}\Bigr) = \frac{31}{90720},
        \qquad
        \frac{\langle Gu, Gu\rangle}{\langle u, Gu\rangle} =
        \frac{31}{306},
        \]
        en besluit $\pi \leq \sqrt{306/31} < 3.1419$: twee
        integralen, vier juiste cijfers. (Elke verdere iteratie
        kwadrateert de nauwkeurigheid ruwweg: de kloof
        $\lambda_1/\lambda_2 = 4$ tussen de eigenvectoren drijft
        een meetkundige convergentie aan.)
\end{enumerate}

\medskip
\noindent\textbf{Deel VI --- Het spoor van $G^2$, en
$\zeta(4)$.}
\begin{enumerate}[resume]
  \item Toon aan dat $\sum_n\lambda_n^2 =
        \iint_{\intcc01^2}g(x,y)^2\,\dd x\,\dd y$ \emph{(ontwikkel
        $g$ op de productbasis $(e_m(x)e_n(y))_{m,n}$ van
        $L^2(\intcc01^2)$ --- een \omterm{def:b3:hilbert:onb}{hilbertbasis}, vgl.\
        \cref{exo:b3:spectral:5} --- en pas \omterm{thm:b3:hilbert:parseval}{Parseval} op het
        vierkant toe; vraag 7 identificeert de coëfficiënten)}.
  \item Bereken de dubbele integraal:
        \[
        \iint g^2 = 2\int_0^1(1-x)^2\Bigl(\int_0^x
        y^2\,\dd y\Bigr)\dd x
        = \frac23\int_0^1x^3(1-x)^2\,\dd x = \frac1{90} .
        \]
  \item Besluit $\zeta(4) = \dfrac{\pi^4}{90}$; leg zonder
        berekening uit hoe de sporen van hogere machten $G^k$
        voor elke $k \geq 1$ $\zeta(2k) \in \pi^{2k}\,\Q$
        voortbrengen, en waarom de oneven waarden $\zeta(3),
        \zeta(5), \dots$ structureel buiten het bereik van deze
        machine liggen.
  \item ($\pi$ van onderen) Leid uit $\lambda_1^2 \leq
        \sum_n\lambda_n^2 = \frac1{90}$ af dat $\pi \geq
        90^{1/4} > 3.080$, en stel met vraag 14 het tweezijdige
        vonnis
        \[
        3.080 \;<\; \pi \;<\; 3.1419
        \]
        samen, \omterm{def:b3:complete:complete}{volledig} verkregen uit de rekenkunde van de
        trillende snaar. Welke kant convergeert sneller als je
        hogere sporen $(\operatorname{tr}G^{2k})^{-1/4k}$
        gebruikt, en waarom?
\end{enumerate}

\medskip
\noindent\textbf{Deel VII --- Aandrijving en resonantie.} Leg
$\nu \in \R$ vast en beschouw de aangedreven snaar $-u'' - \nu u
= f$ met $u(0) = u(1) = 0$, waarbij $f \in L^2$ en $c_n = \langle
e_n, f\rangle$.
\begin{enumerate}[resume]
  \item Stel dat $\nu \notin \{n^2\pi^2 : n \geq 1\}$. Toon aan
        dat
        \[
        u = \sum_{n\geq1}\frac{c_n}{n^2\pi^2 - \nu}\,e_n
        \]
        convergeert in $L^2$ en uniform op $\intcc01$
        \emph{(Cauchy--Schwarz tussen $(c_n)$ en de staarten van
        $\sum n^{-4}$, met $\norm{e_n}_\infty = \sqrt2$)}, en dat
        ze voldoet aan $u = Gf + \nu Gu$ --- de vorm in expliciete
        coördinaten van het \omterm{thm:b3:spectral:fredholm}{alternatief van Fredholm}
        (\cref{thm:b3:spectral:fredholm}), met uniciteit.
  \item Stel dat $\nu = m^2\pi^2$. Toon aan dat $u = Gf + \nu
        Gu$ een oplossing heeft dan en slechts dan als $c_m = 0$,
        uniek op het optellen van veelvouden van $e_m$ na.
        Fysische lezing: een schommel precies op haar eigen
        frequentie aanduwen.
  \item Toon voor $\nu < \pi^2$ aan dat de oplossingsoperator
        $R_\nu\colon f \mapsto u$ begrensd is op $L^2$ met norm
        $\frac1{\pi^2 - \nu}$, en \omterm{def:b3:spectral:compact}{compact}, \omterm{def:b3:spectral:selfadjoint}{zelftoegevoegd} en
        positief: de hele spectraalanalyse begint opnieuw,
        verschoven over $\nu$.
  \item (Synthese) Stel het woordenboek van deze opgave samen:
        eigenwaarde $\leftrightarrow$ gekwadrateerde frequentie
        (harmonischen); spoor $\leftrightarrow$ $\zeta(2)$;
        hilbert-schmidtnorm $\leftrightarrow$ $\zeta(4)$;
        \omterm{thm:b3:spectral:fredholm}{alternatief van Fredholm} $\leftrightarrow$ resonantie;
        min-max $\leftrightarrow$ variationele grenzen ($\pi <
        3.1437$ uit één veelterm). Eén integraaloperator, vijf
        hoofdstukken analyse verzilverd.
\end{enumerate}

\medskip
\noindent\textbf{Deel VIII --- Drie laatste echo's.}
\begin{enumerate}[resume]
  \item ($\zeta(6)$, gratis) De identiteit van \omterm{thm:b3:hilbert:parseval}{Parseval} uit vraag
        6 gaf $\sum_{n\text{ oneven}}n^{-6} = \frac{\pi^6}{960}$.
        Splits $\zeta(6)$ in oneven en even $n$ en besluit
        \[
        \zeta(6) = \frac{\pi^6}{945},
        \]
        zonder één nieuwe integraal: de machine van vraag 17
        (sporen van $G^3$) zou dezelfde waarde hebben opgeleverd,
        tegen de prijs van een geïtereerde kern --- \omterm{thm:b3:hilbert:parseval}{Parseval} op
        één goed gekozen functie is hier de goedkopere weg.
  \item (De grondtoestand is positief) Zij $A$ een \omterm{def:b3:topology:compact}{compacte}
        \omterm{def:b3:spectral:selfadjoint}{zelftoegevoegde} positieve operator op $L^2(\intcc01)$,
        gegeven door een \omterm{def:b3:topology:continuity}{continue} symmetrische kern $k > 0$ op
        $\intoo01^2$, met grootste eigenwaarde $\mu_1$. Toon aan:
        (a) elke maximalisator van het quotiënt van Rayleigh is
        een eigenfunctie bij $\mu_1$; (b) is $u$ er een, dan is
        $\langle\abs u, A\abs u\rangle \geq \langle u, Au\rangle$,
        met \emph{strikte} ongelijkheid als $u$ beide tekens
        aanneemt op verzamelingen van positieve \omterm{def:b3:measure:measure}{maat} --- dus
        heeft $u$ b.o.\ een constant teken, en verdwijnt $u =
        \mu_1^{-1}Au$ nergens op $\intoo01$; (c) $\mu_1$ is een
        \emph{enkelvoudige} eigenwaarde. Ga elke bewering na op
        $G$: $e_1 = \sqrt2\sin(\pi x) > 0$, en elke $e_n$ met $n
        \geq 2$ moet, loodrecht staand op $e_1$, van teken
        wisselen (en dat doet ze: $n - 1$ inwendige nulpunten).
  \item (Afstand tot het spectrum, en de prijs van resonantie)
        Toon voor $\nu \notin \{n^2\pi^2\}$ aan dat de
        oplossingsoperator $R_\nu$ van vraag 19 begrensd,
        \omterm{def:b3:spectral:selfadjoint}{zelftoegevoegd} en \omterm{def:b3:spectral:compact}{compact} is, met
        \[
        \vertiii{R_\nu} =
        \frac1{\min_{n\geq1}\,\abs{n^2\pi^2 - \nu}}
        = \frac1{\operatorname{dist}\bigl(\nu,
        \{n^2\pi^2\}\bigr)},
        \]
        waarbij de norm op de dichtstbijzijnde modus wordt
        bereikt. Kwantificeer dan de schommel van vraag 20:
        aandrijven met $f = e_1$ bij $\nu = (1 - \varepsilon)\pi^2$
        levert $u = \frac{1}{\varepsilon\pi^2}\,e_1$ op, een
        versterking met $\frac1\varepsilon$ van het statische
        antwoord $Ge_1 = \frac1{\pi^2}e_1$ --- één procent onder
        de grondtoon ($\varepsilon = 10^{-2}$) antwoordt de snaar
        honderd keer luider.
\end{enumerate}
\end{problem}
